% !TEX root = ../main.tex
\section{Visão geral e contexto}
\label{sec:S01-visao-geral}

\tool{SAPHO} (\emph{Scalable-Architecture Processor for Hardware Optimization}) é a suíte
integrada para concepção, compilação, simulação e visualização de
\emph{soft-processors} mantida pelo \textbf{NIPSCERN}, núcleo de pesquisa da
\textbf{Universidade Federal de Juiz de Fora (UFJF)}, em Minas Gerais. O
laboratório colabora com o CERN no experimento ATLAS do LHC e emprega a
plataforma no ensino (disciplina DLP --- Dispositivos Lógicos Programáveis). A
expansão da sigla está declarada no campo \code{description} do
\file{package.json}. Versão documentada: \textbf{6.3.2}; licença \textbf{MIT}.

Um \emph{soft-processor} é um processador descrito em HDL (\emph{Verilog}) e
implementado na malha reconfigurável de um FPGA; sua arquitetura (largura, ISA,
periféricos, memória) é parametrizável e recompilável em minutos. Todo o fluxo
\tool{SAPHO} roda localmente, sem nuvem.

\subsection{Modelo conceitual}

\begin{tabularx}{\linewidth}{l Y}
\toprule
\textbf{Peça} & \textbf{Papel} \\
\midrule
\tool{AURORA} & IDE desktop (Electron 39 + editor Monaco; repositório \repo{aurora}) que orquestra escrita, compilação, simulação, ondas e RTL sob interface única. Projetos no tipo \code{.spf} (\emph{SAPHO Project File}), organizados por processador. \\
\tool{YANC} & \enquote{Yet Another Compiler} (em C, Flex/Bison; repositório \repo{yanc}): compila \cmm{} ou C++ até Verilog sintetizável, imagens de memória e \emph{testbench}. Binários: \code{cmmcomp}, \code{asmcomp}, \code{appcomp}, \code{comp2gtkw}. \\
\emph{toolchain bundle} & \emph{Snapshot} portátil msys/mingw64 (\code{aurora-msys-v1.zip}, repositório \repo{aurora-toolchain}) baixado no \emph{bootstrap}: Icarus Verilog, Verilator, Yosys, GCC/G++, \code{make}, Perl, Python+cocotb. \\
\bottomrule
\end{tabularx}

Dois subsistemas operam dentro da IDE. \tool{PRISM} é o visualizador de RTL, que
renderiza a netlist via Yosys + netlistsvg (\emph{fork} \repo{netlistsvg-aurora}).
\tool{Aurora Intelligence} é o subsistema de IA, \textbf{totalmente
implementado} (Vercel AI SDK; provedores \code{openai}, \code{anthropic},
\code{google}, \code{deepseek}, \code{groq}, \code{ollama}; servidor MCP próprio);
porém o \emph{provedor e os modelos próprios} do laboratório ainda \textbf{não
foram treinados} (roadmap), e hoje opera por provedores de terceiros com chaves
do usuário.

\subsection{Fluxo, público e plataforma}

Fluxo de alto nível: \emph{projeto \code{.spf} $\to$ código \cmm{}/C++ $\to$
compilação (YANC) $\to$ Verilog + memórias + \emph{testbench} $\to$ simulação
(iverilog/verilator + cocotb) $\to$ visualização de ondas (GTKWave/Surfer) e RTL
(PRISM)}.

Público-alvo: \textbf{pesquisa e ensino} em arquitetura de computadores e projeto
de soft-processors --- estudantes/docentes e pesquisadores que iteram variações
de arquitetura localmente.

Plataformas: a \textbf{IDE AURORA} hoje roda \textbf{somente em Windows} (alvo
\code{nsis}/\code{x64} no electron-builder); o \tool{YANC} suporta Windows e
Linux. Os repositórios \repo{aurora} (desenvolvimento) e \repo{sapho} (canal de
release) são separados por intenção.
